% CVPR 2025 Paper Template; see https://github.com/cvpr-org/author-kit

\documentclass[10pt,twocolumn,letterpaper]{article}

%%%%%%%%% PAPER TYPE  - PLEASE UPDATE FOR FINAL VERSION
% \usepackage{iccv}              % To produce the CAMERA-READY version
% \usepackage[review]{iccv}      % To produce the REVIEW version
\usepackage[pagenumbers]{cvpr} % To force page numbers, e.g. for an arXiv version

% Import additional packages in the preamble file, before hyperref
%
% --- inline annotations
%
\newcommand{\red}[1]{{\color{red}#1}}
\newcommand{\todo}[1]{{\color{red}#1}}
\newcommand{\TODO}[1]{\textbf{\color{red}[TODO: #1]}}
% --- disable by uncommenting  
% \renewcommand{\TODO}[1]{}
% \renewcommand{\todo}[1]{#1}




% Include other packages here, before hyperref.

% It is strongly recommended to use hyperref, especially for the review version.
% hyperref with option pagebackref eases the reviewers' job.
% Please disable hyperref *only* if you encounter grave issues, e.g. with the
% file validation for the camera-ready version.
%
% If you comment hyperref and then uncomment it, you should delete
% ReviewTempalte.aux before re-running LaTeX.
% (Or just hit 'q' on the first LaTeX run, let it finish, and you
%  should be clear).
\definecolor{cvprblue}{rgb}{0.21,0.49,0.74}
\usepackage[pagebackref,breaklinks,colorlinks,allcolors=cvprblue]{hyperref}
\usepackage{verbatim}
\usepackage{comment}
\usepackage{wrapfig}
\usepackage[accsupp]{axessibility}

\usepackage{graphicx}
\usepackage{amsmath}
\usepackage{amssymb}
\usepackage{booktabs}
% \usepackage{caption}
\usepackage{bm}
\usepackage{adjustbox}
\usepackage{multirow}

\usepackage{url}            % simple URL typesetting
\usepackage{amsfonts}       % blackboard math symbols
\usepackage{nicefrac}       % compact symbols for 1/2, etc.
\usepackage{microtype}      % microtypography
\usepackage{xcolor}         % colors

\usepackage{array}
\usepackage{xspace}


% Support for easy cross-referencing
% \usepackage[capitalize]{cleveref}
% \crefname{section}{Sec.}{Secs.}
% \Crefname{section}{Section}{Sections}
% \Crefname{table}{Table}{Tables}
% \crefname{table}{Tab.}{Tabs.}


\usepackage[ruled]{algorithm2e} % For algorithms



% require xspace, array

\newcommand{\stderr}[1]{\scriptsize $\pm #1$}

%% layout
\newcolumntype{L}[1]{>{\raggedright\let\newline\\\arraybackslash\hspace{0pt}}m{#1}}
\newcolumntype{C}[1]{>{\centering\let\newline\\\arraybackslash\hspace{0pt}}m{#1}}
\newcolumntype{R}[1]{>{\raggedleft\let\newline\\\arraybackslash\hspace{0pt}}m{#1}}

\newcommand{\xpar}[1]{\noindent\textbf{#1}\ \ }
\newcommand{\vpar}[1]{\vspace{3mm}\noindent\textbf{#1}\ \ }

%% notations
\newcommand{\sect}[1]{Section~\ref{#1}}
\newcommand{\sectapp}[1]{Appendix~\ref{#1}}
\newcommand{\ssect}[1]{\S~\ref{#1}}
\newcommand{\eqn}[1]{Equation~\ref{#1}}
\newcommand{\fig}[1]{Figure~\ref{#1}}
\newcommand{\tbl}[1]{Table~\ref{#1}}
\newcommand{\alg}[1]{Algorithm.~\ref{#1}}

\newcommand{\degree}{\ensuremath{^\circ}\xspace}
\newcommand{\ignore}[1]{}
\newcommand{\norm}[1]{\lVert#1\rVert}
\newcommand{\fcseven}{$\mbox{fc}_7$}
% \DeclarePairedDelimiter\ceil{\lceil}{\rceil}
% \DeclarePairedDelimiter\floor{\lfloor}{\rfloor}

\renewcommand*{\thefootnote}{\fnsymbol{footnote}}

%\DeclareMathOperator*{\argmin}{arg\,min}
%\DeclareMathOperator*{\argmax}{arg\,max}

\def\naive{na\"{\i}ve\xspace}
\def\Naive{Na\"{\i}ve\xspace}

\makeatletter
\DeclareRobustCommand\onedot{\futurelet\@let@token\@onedot}
\def\@onedot{\ifx\@let@token.\else.\null\fi\xspace}

\def\eg{e.g\onedot} \def\Eg{E.g\onedot}
\def\ie{i.e\onedot} \def\Ie{I.e\onedot}
\def\cf{\emph{c.f}\onedot} \def\Cf{\emph{C.f}\onedot}
\def\etc{etc\onedot}
\def\vs{\emph{vs}\onedot}
\def\wrt{w.r.t\onedot}
\def\dof{d.o.f\onedot}
\def\etal{et al\onedot}
\def\aka{a.k.a\onedot}
\makeatother

%% comments
\definecolor{MyDarkBlue}{rgb}{0,0.08,1}
\definecolor{MyDarkGreen}{rgb}{0.02,0.6,0.02}
\definecolor{MyDarkRed}{rgb}{0.8,0.02,0.02}
\definecolor{MyDarkOrange}{rgb}{0.40,0.2,0.02}
\definecolor{MyPurple}{RGB}{111,0,255}
\definecolor{MyRed}{rgb}{1.0,0.0,0.0}
\definecolor{MyGold}{rgb}{0.75,0.6,0.12}
\definecolor{MyDarkgray}{rgb}{0.66, 0.66, 0.66}

\definecolor{turquoise}{cmyk}{0.65,0,0.1,0.3}

\newcommand{\karen}[1]{\textcolor{red}{[Karen: {#1}]}}
\newcommand{\jiaman}[1]{\textcolor{orange}{[Jiaman: {#1}]}}
\newcommand{\zhen}[1]{\textcolor{blue}{[Zhen: {#1}]}}

\renewcommand{\algorithmcfname}{ALGORITHM}

%%%%%%%%% PAPER ID  - PLEASE UPDATE
\def\paperID{6622} % *** Enter the CVPR Paper ID here
\def\confName{ICCV}
\def\confYear{2025}


% Document starts

% Title portion
\title{Human-Object Interaction from Human-Level Instructions}

% \author{First Author\\
% Institution1\\
% Institution1 address\\
% {\tt\small firstauthor@i1.org}
% % For a paper whose authors are all at the same institution,
% % omit the following lines up until the closing ``}''.
% % Additional authors and addresses can be added with ``\and'',
% % just like the second author.
% % To save space, use either the email address or home page, not both
% \and
% Second Author\\
% Institution2\\
% First line of institution2 address\\
% {\tt\small secondauthor@i2.org}
% }

\author{Zhen Wu \qquad Jiaman Li\qquad Pei Xu\qquad C. Karen Liu \\
Stanford University\\
% Institution1 address\\
{\tt\small \{zhenwu,jiamanli,peixu,karenliu\}@cs.stanford.edu}}

% \begin{document}

% \twocolumn[{%
% \renewcommand\twocolumn[1][]{#1}%
% \maketitle
% \begin{center}
%     \centering
% \includegraphics[width=\linewidth]{figures/0_teaser.pdf}
% % \vspace{-2em}
% \captionof{figure}{\zhen{need revision}}
% \label{fig:teaser}
% \end{center}
% }]

\begin{document}
% \maketitle
\twocolumn[{%
\renewcommand\twocolumn[1][]{#1}%
% \vspace{-3em}
\maketitle
\vspace{-12mm}
\begin{center}
    \centering
    \includegraphics[width=\textwidth]{figures/0_teaser_v1_compressed.pdf}
    \vspace{-6mm}
    \captionof{figure}{Given human-level instructions (top), we first use a diffusion-based motion generator to synthesize natural full-body human motion, finger motion, and object motion to accomplish the task (middle). Next, we employ a physics-based tracking method to imitate the generated interactions which produces physically plausible interactions (bottom). Please refer to our \href{https://hoifhli.github.io/}{\textcolor{magenta}{project page}} for more results.}
    \label{figure1}
    \vspace{-6pt}
\end{center}%
}]

\setcounter{figure}{1}

\begin{abstract}
Intelligent agents must autonomously interact with the environments to perform daily tasks based on human-level instructions. They need a foundational understanding of the world to accurately interpret these instructions, along with precise low-level movement and interaction skills to execute the derived actions.
In this work, we propose the first complete system for synthesizing physically plausible, long-horizon human-object interactions for object manipulation in contextual environments, driven by human-level instructions. We leverage large language models (LLMs) to interpret the input instructions into detailed execution plans. Unlike prior work, our system is capable of generating detailed finger-object interactions, in seamless coordination with full-body movements. We also train a policy to track generated motions in physics simulation via reinforcement learning (RL) to ensure physical plausibility of the motion. Our experiments demonstrate the effectiveness of our system in synthesizing realistic interactions with diverse objects in complex environments, highlighting its potential for real-world applications.
\vspace{-2mm}

\end{abstract}



\section{Introduction}
\label{sec:introduction}

% (1 Background) 
Synthesizing physically plausible, long-horizon, human-object interactions in contextual environments is crucial for applications in computer graphics, embodied AI, and robotics. When given human-level language instructions, intelligent agents should be able to comprehend these high-level commands and map them to executable action sequences. For instance, when instructed to set up a workspace, an agent must understand the abstract concept of ``workspace'', identify the relevant furniture and objects in the room, and translate this concept into a series of concrete actions, such as moving and orienting a desk, placing a chair nearby, and positioning a monitor on the desk. In addition to understanding the instructions, the agent must execute actions that lead to realistic and functional human movements, including navigating cluttered spaces, manipulating various objects for the task, and appropriately releasing them after use.
% (2 What prior work do and cannot achieve) 

Prior research has made significant progress in synthesizing human-scene interactions using large language model (LLM) planners~\cite{xiao2024unified}. However, most of these interactions have been limited to static objects without hand manipulation. Some studies~\cite{li2023controllable} have demonstrated the manipulation of dynamic objects, but the resulting interactions still lack realistic finger movements. Recently, physics-based RL policies has demonstrated promising results in human-object interactions~\cite{luo2024grasping, wang2023physhoi}. However, generating long-horizon motions that interact with various objects in a scene-aware manner remains an ongoing challenge. Build upon these prior work, we introduce the first complete system capable of synthesizing physically accurate, long-horizon human-object interactions with synchronous full-body and intricate finger movements, from human-level language instructions. 



Synthesizing coordinated full-body and finger movements from human-level instructions presents two main challenges. First, human-level instructions typically provide only a high-level task outline, requiring common sense and world knowledge to translate them into precise object arrangements and detailed execution plans. Second, existing datasets with paired full-body and finger motion data are limited to interactions with small objects and lack locomotion during object manipulation~\cite{GRAB:2020, fan2023arctic}.
Although a recent dataset~\cite{li2023object} includes locomotion and manipulation for large objects, it lacks detailed finger movements due to the difficulty of capturing both modalities simultaneously. 

To tackle the challenge of comprehending human-level instructions, we utilize the latest advancements in LLMs, which are highly effective in interpreting high-level commands\cite{openai2023gpt, liu2024visual}.
However, LLMs are not well-suited for directly predicting precise scene layouts, such as the exact positions and orientations of objects~\cite{aguina2024open, hong20233d}.
Instead of prompting LLMs to directly specify object placements, we use object spatial relationships as intermediate representations and ask LLMs to derive these relationships. We then propose an algorithm to calculate the precise object poses based on these relationships.

For the low-level motion synthesis, we take the approach of modeling human motion from real-world human data. To address the lack of large-scale object manipulation datasets with synchronized full-body and finger motions, one could use two separate datasets to train a full-body manipulation model and another detailed finger manipulation model. 
% However, a causality dilemma emerges when putting the motions together---should we first generate the full-body manipulation and then enforce the finger-object contacts to agree with the wrist poses determined by the full-body motion, or should we proceed with the other way around? In either approach, simply fusing the two motions together will cause glaring artifacts due to the disagreement in two separately trained models.
However, simply fusing the two motions together will cause glaring artifacts due to the disagreement in two separately trained models.
%The former would result in glaring physical artifacts at finger-object contacts while the latter requires both datasets to share the same object motion. 
% Our solution to break the dilemma is to employ \emph{both} approaches in three steps. 
Our solution is to break this process into three steps.
We first generate full-body and object motions without detailed fingers. The object motion and rough wrist poses provide guidance for generating contact-consistent finger motion in the second step. Finally, we generate full-body and object motions again conditioned on precise finger motions. This approach enables our framework to successfully produce synchronized object motion, full-body motion, and finger motions without requiring a dataset that includes all data modalities.

% However, existing datasets only consist of kinematic pose trajectories, resulting in generative models that do not produce physically accurate motions. 
The generated motion, however, is not guaranteed to be physically accurate.
To ensure physical plausibility, we further use reinforcement learning (RL) to track the generated motion in a physics simulator. 
Our method adopts an importance sampling technique to improve the training efficiency of long sequence tracking and supports tracking diverse interactions with multiple distinct objects within a single sequence.

We evaluate our method through extensive quantitative and qualitative comparisons, as well as human studies, demonstrating that our approach outperforms all others. This underscores the effectiveness of our method and its potential for real-world applications.

% (5 Summary of Contributions) 
%To summarize, our work makes the following contributions: (i) We leverage LLMs to translate human-level instructions into precise spatial task planning grounded in executable actions. 
%(ii) We propose an effective motion generator that synthesizes human-object interactions with detailed finger movements. (iii) We apply RL to ensure the physical plausibility of the generated motions. (iv) By integrating these components, we develop the first complete system capable of synthesizing physically accurate, long-horizon human-object interactions with intricate finger movements, from human-level language instructions. \karen{We can just talk about evaluation in this paragraph.}

% \vspace{10pt}

\section{Related Work}
\label{sec:related_work}

\paragraph{Contextual Interaction Synthesis.}
% \paragraph{Scene-Aware Human Motion Synthesis. }
% Based on paired scene and human motion datasets~\cite{}, synthesizing human motions in 3D scenes has been extensively explored~\cite{} in recent years. However, the paired datasets were captured by filming multi-view RGBD videos and conducting optimization to extract human motions, which led to inferior motion quality. To improve the motion quality, prior work explored combining 3D scene datasets~\cite{} and motion capture dataset~\cite{} to synthesize paired data. A novel combination of VR headset and motion capture was explored in prior work~\cite{}. Recent work scaled up the data collection for paired scene and human motion by building physical scenes in the mocap lab to obtain high-quality interaction motion. 
Recent work on modeling human-object interactions falls into two categories: interactions with static and dynamic objects. Synthesizing human motions in static 3D scenes has been extensively studied~\cite{prox, wang2022humanise, araujo2023circle, jiang2024scaling, hassan_samp_2021, zhang2022couch}. 
Prior research has explored regression models~\cite{wang2021synthesizing, wang2022towards, hassan_samp_2021, araujo2023circle, mir2023generating} and diffusion models~\cite{huang2023diffusion, jiang2024scaling, yi2024generating, wang2024move} to generate human motions interacting with 3D scenes, such as sitting on a chair or lying down on a sofa. Some studies have also explored reinforcement learning methods to generate physically plausible human motions~\cite{hassan2023synthesizing, zhao2023synthesizing, xiao2024unified}. These works focus on static objects, whereas our work focuses on dynamic interactions.


There has been increasing attention on synthesizing human-object interactions with dynamic objects~\cite{diller2023cg, peng2023hoi, wu2024thor,xu2024interdreamer}. OMOMO~\cite{li2023object} collected a dataset for dynamic human-object interactions and proposed a framework to synthesize human motions from object motions. 
% InterDiff~\cite{xu2023interdiff} predicts future interaction motions based on past object and human motion.
CHOIS~\cite{li2023controllable} synthesizes interactions in 3D scenes from text and sparse waypoints.
Despite these advancements, all of these works focus on full-body human motion generation without detailed finger motions. In this work, we aim to generate synchronized object motions, human motions, and finger motions.

\paragraph{Hand-Object Interaction Synthesis. }
Hand-object interaction synthesis has been extensively explored.
In the realm of static pose generation, traditional methods often rely on physics-based control or optimization techniques to generate grasp candidates \cite{li2007data, pollard2005physically, rodriguez2012caging}. 
% Some recent work utilizes differentiable force closure estimators to generate high-quality grasp pose \cite{wang2023dexgraspnet, liu2021synthesizing}.
Some recent work leverages deep learning to directly estimate hand grasps from extensive interaction datasets \cite{jiang2021hand, brahmbhatt2019contactgrasp, karunratanakul2020grasping, GRAB:2020, fan2023arctic, christen2024diffh2o}. 
Beyond static grasping, another line of work also explores dynamic object manipulation \cite{zhang2021manipnet, ye2012synthesis, liu2024geneoh}. 
% ManipNet \cite{zhang2021manipnet} introduces a spatial representation based on the hand-object distance and trains an auto-regressive model to generate finger movement. 

Recent works further explore integrating whole-body motion and hand motion \cite{taheri2023grip, tendulkar2023flex, wu2022saga, taheri2022goal}.
IMoS \cite{ghosh2023imos} generates hand-object interaction using paired human motion and hand motion, while~\cite{braun2023physically} leverage reinforcement learning to achieve physically plausible motion.
% However, these methods are typically trained on paired data of whole-body and hand motion, making them challenging to generalize to different scenarios, such as interacting with large objects.
However, these methods are constrained to manipulate small-sized objects in the GRAB dataset~\cite{GRAB:2020}, unable to generalize to large objects. 
In this work, we aim to generate realistic full-body and finger motions for manipulating large objects. Since no such paired dataset exists, we leverage DexGraspNet~\cite{wang2023dexgraspnet} to generate realistic grasp poses, train a conditional diffusion model on GRAB~\cite{GRAB:2020} to generate approaching and releasing finger motions, and then integrate these with our full-body motion generation.   
% FLEX \cite{tendulkar2023flex} proposes a method to separate hand motion and whole-body motion generation. It utilizes distinct datasets to train models for each, which are then combined to generate single-frame full-body grasp poses.
% \zhen{needs to revise: 
% In this work, we extend this idea to motion sequence generation. We utilize hand-object interaction datasets such as GRAB \cite{taheri2020grab} and ARCTIC \cite{fan2022articulated} to train a finger motion generator, which is then combined with a whole-body motion generator. Additionally, we introduce a novel grasp-guided finger motion generation method, significantly improving the quality of generated hand motion.}

% \paragraph{Diffusion-Based Motion Synthesis.}
% Diffusion models have shown strong capabilities in human motion synthesis~\cite{shafir2023human, tevet2022human, zhang2022motiondiffuse, zhang2023tedi, shi2023controllable, karunratanakul2023gmd, chen2023executing}, with several models addressing interactions with objects or scenes~\cite{li2023controllable, li2023object, xu2024interdreamer}. 
% TeSMo~\cite{yi2024generating} leverages a pre-trained motion diffusion model and finetunes a scene-aware branch to incorporate both text control and scene awareness.
% \cite{jiang2024autonomous} employs an auto-regressive diffusion model to synthesize the next motion segment in a 3D scene.
% In this work, we use the diffusion model as a key component of our low-level motion generator. \jiaman{Do we need this paragraph? We already talked about some diffusion-based interaction synthesis in the first paragraph.} \zhen{A siggraph asia reviewer wanted to see that. What do you all think?}

\paragraph{Physics-Based Motion Synthesis.}
Physics-based methods control a simulated character to perform motion in a simulator, ensuring the physical realism of the generated motion. To make the motion more natural, many approaches leverage motion capture data to produce movements that match the style of the data~\cite{peng2018deepmimic, peng2021amp, peng2023hoi, liu2018learning}. DeepMimic~\cite{peng2018deepmimic} proposes to track motion capture data using RL. Subsequent works~\cite{peng2021amp, xu2021gan} leverage generative adversarial networks (GANs) to learn motion styles. Recent research also explores this approach for human-object interactions, such as playing musical instruments~\cite{xu2024synchronize, wang2024f} or playing soccer~\cite{xie2022learning}.  PhysHOI~\cite{wang2023physhoi} learns basketball-playing skills from human demonstrations. OmniGrasp~\cite{luo2024grasping} enables a character to grasp and move an object along a specified object trajectory.
In this work, we leverage RL to track generated kinematic motion, successfully producing long-horizon motions that interact with various objects.

\paragraph{LLM-based Planning.}
LLMs have been widely used in reasoning, planning tasks~\cite{hong20233d, hu2023look, wake2023gpt, zeng2023distilling, ding2023task} for their high-level planning capabilities.
They have also been adopted in motion synthesis tasks~\cite{yao2024moconvq, zhang2024motiongpt, xiao2024unified, xu2024interdreamer}.
UniHSI~\cite{xiao2024unified} uses LLMs to extract action plans in a chain-of-contact format, guiding a low-level controller to synthesize corresponding motions. 
InterDreamer~\cite{xu2024interdreamer} employs LLMs to identify object parts that humans interact with and uses this information to retrieve initial human poses. 
% CHOIS~\cite{li2023controllable} uses LLMs to extract key information to determine target object position, deriving sparse waypoints as input for a conditional diffusion model to generate interactions grounded in scenes. However, CHOIS requires explicit language descriptions (e.g., "Lift the monitor, move, and put it on the desk.") as input, which cannot support interaction synthesis from more abstract human-level instructions (e.g., "Set up a workspace"). 
In this work, we leverage LLMs to translate abstract human-level instructions into target scene maps and detailed execution plans, which guide the low-level interaction synthesis.


\section{System Overview}
\label{sec:system_overview}

\begin{figure*}[t]
    \centering
    \includegraphics[width=\textwidth]{figures/1_system_overview_v7_compressed.pdf}
    \vspace{-6mm}
    \caption{
    Our system takes the scene description and human-level instruction as input and uses a high-level planner to obtain the scene map and a detailed execution plan. The low-level motion generator then generates synchronized object motion, full-body human motion, and finger motion. Finally, the physics tracker uses RL to track the generated motion, producing physically plausible motion.}
    % \Description{}
    \vspace{-4mm}
    \label{fig:system_overview}
\end{figure*}
Given a human-level instruction and a scene description, we develop a system that generates synchronized human motion and object motion to accomplish the task. 
%Our system differs from prior work in human-scene and human-object interaction by not only modeling the full-body dynamics of human motion but also focusing on the intricate details of finger movements. This enhancement contributes significantly to the realism of the motions generated.
An overview of our system is depicted in \cref{fig:system_overview}. The system consists of a high-level planner, a low-level motion generator, and a physics tracker. The high-level planner leverages LLMs to reason about instructions and the scene description, generating a \emph{scene map} along with a detailed \emph{execution plan}. The scene map specifies the desired positions and orientations of the objects, while the execution plan outlines the steps to interact with these objects. The low-level motion generator unites conditional diffusion models trained on separate datasets for full body motion and dexterous hand motion, producing synchronous object and human movements.
%It includes an interaction module and a navigation module. The interaction module, introduced in \cref{sec:low_level_interaction}, employs a multi-stage framework to produce human-object interaction motions. The navigation module, described in \cref{sec:low_level_navigation}, generates walking motions between interactions.
The physics tracker then uses RL to track the generated motion in a physics simulator, ensuring physical plausibility.
% Together, these modules operate sequentially to ensure smooth and realistic human-object interactions.







\section{High-Level LLM Planner}
\label{sec:high_level_planner}


Inputs to the high-level planner consist of an initial scene description and human-level instructions. The scene description includes a list of objects in a 3D scene and their layout, while the human-level instruction describes a task that requires grounded reasoning based on world knowledge and common sense. For example, the instruction ``I want to do yoga in front of the TV'' implies that any obstacles in front of the TV should be moved elsewhere.

Outputs of the high-level planner comprise the scene map and the execution plan.
The scene map specifies the target position $p$ and orientation $q$ for each object $o$ that requires movement, represented as $\left\{ \left\{o_1, p_1, q_1 \right\}, \cdots, \left\{o_n, p_n, q_n \right\}\right\}$.
The execution plan is a sequence of text actions associated with objects in the scene map. For example, “lift the monitor, move the monitor, put down the monitor”.

\subsection{Scene Map Generation}
\label{subsec:spatial_relationship}

% Given a human-level instruction and a scene description, ideally, LLMs should directly output object positions and orientations.
% However, recent research suggests that state-of-the-art LLM's ability to reason 3D information is not yet reliable \cite{aguina2024open, hong20233d}.
% To increase LLM's reliability in spatial reasoning for our tasks, we use a in-context learning strategy and propose a method to ensure LLM's answers are actionable by the low-level motion generator. 
To obtain a reliable scene map, inspired by recent works in room layout generation \cite{aguina2024open, yang2024holodeck}, we first instruct LLMs to reason about spatial relationships of the objects as an intermediate representation from which the 3D positions and orientations of the objects are calculated subsequently.
We define the following relation functions to capture spatial relationships, with flexibility to include additional relationships if needed.
% We define the following relation functions to describe the spatial relationships which are sufficient to address most rearrangement tasks in a typical home environment, though our framework remains flexible to include additional relationships if necessary.
% \begin{enumerate}
%     \item \texttt{on(object1, object2)} indicates that \texttt{object1} is positioned on top of \texttt{object2}.
%     \item \texttt{adjacent(object1, object2, direction, distance)} determines the adjacency between \texttt{object1} and \texttt{object2}, with optional parameters, \texttt{direction} and \texttt{distance}, to specify the relative orientation and separation.
%     \item \texttt{facing(object1, object2)} denotes that \texttt{object1} is oriented towards \texttt{object2}.
% \end{enumerate}

\begin{enumerate}
    \item \texttt{on(object1, object2)}.
    \item \texttt{adjacent(object1, object2, direction, distance)}.
    \item \texttt{facing(object1, object2)}.
\end{enumerate}

For example, if a monitor needs to be put on the table facing the chair, we can use \texttt{on(monitor, table)} and \texttt{facing(monitor, chair)} to describe its pose. If a table is positioned 1 meter to the north of the door, we can use \texttt{adjacent(table, door, north, 1)} to describe the relationship. 
We prompt LLMs to generate the spatial relationships of the objects according to the given instruction. 
% In-context learning is employed during the process to assist LLMs to better understand the output requirements. We provide complete prompts for the LLM planner in the supplementary materials. 

% \subsection{Target Layout Calculation}
% \label{subsec:layout}


Given the spatial relationships, we propose an algorithm to calculate the 3D positions of each object in the scene. The orientations of objects, enforced by ``facing'' relations, will be addressed in the subsequent step. 

The algorithm first constructs a scene graph \cite{johnson2015image} to organize the objects (nodes) and their spatial relationships (edges), which are directional, pointing from \texttt{object2} to \texttt{object1}. The algorithm employs a method similar to topological sorting, assuming no cycles are present in the graph. The scene graph is initialized by visiting the nodes associated with static objects and setting their 3D positions as the given values. In each iteration, the algorithm looks for a new node whose 3D position has not been determined but all its predecessors have. By calculating an offset from the predecessors' known positions, we can determine the position of the new node. For example, if object $o_1$ is ``on'' object $o_2$, then $o_1$'s height is set to match the top surface of $o_2$, while the horizontal position is sampled within the horizontal range of $o_2$. If $o_1$ is ``adjacent'' to $o_2$, we use the direction and distance given by the LLMs as the offset. We have included pseudo-code in the supplementary material.
%A post-process is applied to ensure no overlap between objects.

To compute the orientation enforced by \texttt{facing($o_1$, $o_2$)}, we rotate the canonical direction of $o_1$ in its own local frame to align with the vector from the position of $o_1$ to the position of $o_2$. The canonical direction of each object is part of the input 3D geometry information.


\subsection{Execution Plan Generation}
\label{subsec:plan}
 
After determining the scene map, the planner must specify the interaction order with each object to ensure natural motion. For instance, if a vase is on a table, the vase should be moved before the table. 
LLMs are asked to generate this action sequence in a natural order and provide reasoning for it.
The resulting execution plan consists of textual actions, $\left\{l_1, l_2, \dots, l_{T}\right\}$, where each action $l_t$ defines interaction involving one object. 
Each $l_t$ serves as the textual input for the low-level generator, following the format: ``lift the \textit{object}, move the \textit{object}, put down the \textit{object}''. This format is designed to closely match the textual style of the training dataset of the low-level generator.

With the scene map and execution plan in place, an A* planner generates collision-free paths to execute the plan. The path is a series of 2D waypoints in the scene, which then serve as inputs for the low-level motion generator.
\section{Low-Level Motion Generator}

\begin{figure*}[t]
    \centering
\includegraphics[width=\textwidth]{figures/3_coarse2fine_v7_compressed.pdf}
\vspace{-6mm}
    \caption{The interaction module consists of four stages. Starting with the initial human and object poses, waypoints, final object pose, and a text instruction, the first stage generates coarse object and human motion. The relative wrist pose from this motion is then used to generate a valid grasp pose. Using this as a condition, the next stage re-generates the motion to align with the grasp constraint. Finally, the finger motion generation module produces finger motions during the approach and release of objects.} 
    \vspace{-4mm}
    \label{fig:coarse_to_fine}
\end{figure*}

Given the waypoints and textual instructions generated by the high-level planner, the low-level motion generator executes an interaction module and a navigation module sequentially to generate long-horizon human-object interaction motions.

\subsection{Background: Conditional Diffusion Model}
\label{sec:cond_diff}
The diffusion model involves a forward diffusion process and a reverse diffusion process. Starting from clean data $\bm{x}_0$,  noise is progressively introduced until, after N steps, the data becomes approximately $\bm{x}_{N} \sim \mathcal{N}(0, \bm{I})$. The denoising network is trained to reverse this process so that sampling a noise $\bm{x}_{N}$ and applying N denoising steps yields the clean data representation $\hat{\bm{x}}_0$. Each denoising step is defined as 
\begin{equation}
p_{\theta}(\bm{x}_{n-1}|\bm{x}_{n}, \bm{c}) := \mathcal{N}(\bm{x}_{n-1}; \bm{\mu}_{\theta}(\bm{x}_n, n, \bm{c}), \sigma_{n}^{2}\bm{I}),  
\end{equation}
where $p_{\theta}$ represents a neural network, $\bm{c}$ represents the conditions, $\bm{\mu}_{\theta}(\bm{x}_n, n, \bm{c})$ represents the learned mean, $\sigma_{n}^{2}$ denotes a fixed variance. The training loss is defined as $\mathcal{L} =  \mathbb{E}_{\bm{x}_0, n}||\hat{\bm{x}}_{\theta}(\bm{x}_{n}, n, \bm{c}) - \bm{x}_{0}||_{1}$.
% The detailed architecture is provided in the supplementary material. 

\subsection{Interaction Module}
\label{sec:low_level_interaction}
We introduce a multi-stage interaction module that synthesizes synchronized full-body, finger, and object motion, as shown in~\cref{fig:coarse_to_fine}.
The first stage, \textit{CoarseNet}, generates initial human and object motions without detailed finger movement.
Recognizing that finger poses typically remain static during interaction, the second stage generates a \textit{grasp} pose based on the initial results.
Then \textit{RefineNet} re-generates the human and object motions to seamlessly align with the optimized grasp pose.
Finally, \textit{FingerNet} generates smooth finger motions for approaching and releasing the object, completing the interaction sequence.

\subsubsection{CoarseNet}
\label{sec:coarse_net}
Given the initial human and object poses, waypoints, and final object pose, along with a text instruction, CoarseNet generates synchronized human motion and object motion. This is achieved using a pre-trained diffusion model CHOIS~\cite{li2023controllable}.
% We utilize the pre-trained CHOIS model to generate synchronized object and human motions based on the initial state, end object state, sparse waypoints, object geometry, and text.

% \paragraph{Motion Representation}
% We denote human motion as $\bm{H} \in \mathbb{R}^{T \times D}$, where $T$ represents the sequence length and $D$ represents the dimension of the human pose. Each $\bm{H}_{t}$ includes global joint positions and 6D continuous rotations~\cite{zhou2019continuity}, excluding finger joints. Object motion $\bm{O} \in \mathbb{R}^{T \times 12}$ includes the object's global position and relative rotation matrix with respect to the input object geometry. The model also predicts contact labels $\bm{L} \in \mathbb{R}^2$ for both hands, which are used for test-time guidance. 

% \paragraph{Object Geometry Representation. } Object geometry $\bm{G} \in \mathbb{R}^{1024 \times 3}$ is represented using the Basis Point Set (BPS)~\cite{prokudin2019efficient} method. This method involves sampling $1024$ basis points from the volume of a ball with a $1$-meter radius centered on the object. Directional vectors are then calculated from each basis point to the nearest point on the object, forming a compact BPS representation. 

% \paragraph{Input Condition Representation. } 
% The model uses a masked motion data representation $\bm{S} \in \mathbb{R}^{T \times (D+12)}$ to encode the initial states, end object state, and sparse waypoints. The initial state includes the human pose and object pose in the first frame, while the sparse waypoints provide a series of 2D object positions every 30 frames. Any remaining spaces in $\bm{S}$ are padded with zeros. The masked motion data $\bm{S}$, the BPS representation $\bm{G}$, and text embeddings extracted by a pretrained CLIP model~\cite{radford2021learning} together form the condition $\bm{c}$ for the conditional diffusion model.

% The conditional diffusion model uses input conditions $\bm{c}$ to generate object motion $\bm{O}$, human motion $\bm{H}$, and hand contact $\bm{L}$. To ensure realistic hand-object contact and reduce artifacts, CHOIS employs test-time guidance~\cite{ho2022video} to minimize wrist-object distance in the contact phase, feet-floor distance, and penalize object-floor penetration. For more details, please refer to~\cite{li2023controllable}.


\paragraph{Data Representation.} 
We denote human motion as $\bm{H} \in \mathbb{R}^{T \times D}$, where $T$ represents the sequence length and $D$ represents the dimension of the human pose. Each $\bm{H}_{t}$ includes global joint positions and 6D rotations~\cite{zhou2019continuity}, excluding finger joints. Object motion $\bm{O} \in \mathbb{R}^{T \times 12}$ includes the object's global position and rotation matrix.
CoarseNet predicts human motion, object motion, and contact labels $\bm{L} \in \mathbb{R}^2$ for both hands, resulting in $\bm{x} = \left\{\bm{H}, \bm{O}, \bm{L}\right\}$.

The input condition is represented as $\bm{c} = \left\{\bm{S}, \bm{G}, \bm{T}\right\}$. $\bm{S} \in \mathbb{R}^{T \times (D+12)}$ is a masked motion data representation. It includes the human and object pose in the first frame, the object pose in the last frame, and 2D waypoints sampled every 30 frames. Any remaining entries in $\bm{S}$ are padded with zeros. The object geometry $\bm{G}$ represents the BPS representation~\cite{prokudin2019efficient} to encode object geometry, and $\bm{T}$ refers to the text embeddings extracted using CLIP~\cite{radford2021learning}.

\vspace{-3mm}
\paragraph{Segmenting Contact Phases.} 
% Since the object is stationary at the beginning and end, 
We segment the generated interaction sequence into \textit{pre-contact}, \textit{contact}, and \textit{post-contact} phases based on the predicted contact labels $\bm{L}$.  
% The contact phase is computed separately for each hand. 
% When both hands contact the object, the contact phase is the union of left-hand and right-hand contact phases.
The contact phase is determined separately for each hand, dividing the model output $\bm{x}_\text{coarse}$ into $\bm{x}_\text{pre-contact} \oplus \bm{x}_\text{contact} \oplus \bm{x}_\text{post-contact}$.
% From $\bm{x}_\text{contact}$, we calculate the average wrist pose in the object frame, $\bm{w}$, which is then used as the input for the second stage.
We then compute the average wrist pose $\bm{w}$ from $\bm{x}_\text{contact}$, serving as input for the next stage. 

\subsubsection{Grasp Pose Generation}
\label{subsubsec:grasp_pose}
% The output of CoarseNet provides coordinated full-body and object motion, but the hands maintain a rest pose throughout the motion, resulting in visually and physically implausible hand-object contacts. We address this issue in the second stage by applying a state-of-the-art grasp pose generation method~\cite{wang2023dexgraspnet} to obtain a detailed grasp pose with accurate finger-object contact.

% Based on the object's velocity in the first stage, we segment the generated interaction sequence into pre-contact, contact, and post-contact phases, where the object velocities in the pre-contact and post-contact phases are approximately zero. 
% We also use the average relative wrist pose computed from Stage 1 to generate a plausible grasp pose as an initial guess to be optimized.
% Starting with an object mesh and the initial hand pose $\bm{g} = (\bm{T}, \bm{R}, \bm{\theta})$, where $\bm{T} \in \mathbb{R}^3$ and $\bm{R} \in SO(3)$ represent the wrist position and rotation in the object's frame, and $\bm{\theta} \in \mathbb{R}^{45}$ represents the finger configuration, we formulate an optimization to compute a hand pose that minimize a differentiable force closure energy term and a regularization term (details in \cite{wang2023dexgraspnet}). For motions involving both hands, such as carrying a box, we omit the force closure term because the grasp stability does not rely solely on the fingers of one hand. The optimized grasp pose is denoted by $\hat{\bm{g}} = (\bm{\hat{T}}, \bm{\hat{R}}, \bm{\hat{\theta}})$. The optimized wrist pose $\hat{\bm{w}} = (\bm{\hat{T}}, \bm{\hat{R}})$ is subsequently used as the input condition for the third stage.

The second stage generates the finger pose for the contact phase using a state-of-the-art grasp pose generation method~\cite{wang2023dexgraspnet}.
% Although originally designed for robotic hand grasping, it relies solely on kinematic chain information, making it compatible with SMPL-X~\cite{smplx}.
Given an object, wrist pose $\bm{w}$ and rest finger pose, an optimization is conducted to generate a physically plausible grasp pose.  
We modified \cite{wang2023dexgraspnet} to fit our needs by: (1) increasing the object surface sample points from 2,000 to 20,000 to maintain sample density for larger objects and (2) omitting the force closure term for tasks involving both hands, such as carrying a box, where stability doesn’t solely rely on a single hand.
% We made the following modifications to adapt~\cite{wang2023dexgraspnet} for our task:
% We made several modifications to tailor the optimization method to our specific application:
% \begin{itemize} 
% \item DexGraspNet~\cite{wang2023dexgraspnet} randomly samples 2,000 points on the object surface to measure hand-object penetration. Since our objects are much larger, we increase this to 20,000 sample points to maintain a similar sample density. 
% \item DexGraspNet~\cite{wang2023dexgraspnet} uses a force closure term in the optimization process to measure grasp stability. However, our tasks involve motions with both hands, such as carrying a box, where stability is not solely dependent on one hand. We omit this term for such motions.
% \end{itemize}

The optimized grasp pose is denoted by $\hat{\bm{g}} = (\hat{\bm{w}}, {\hat{\bm{\theta}}})$, where $\hat{\bm{w}}$ is the optimized wrist pose, and $\hat{\bm{\theta}}$ is the finger pose. This grasp pose $\hat{\bm{g}}$ is maintained throughout the contact phase. However, as CoarseNet’s generated motion may not fully align with $\hat{\bm{w}}$, we introduce RefineNet to address this misalignment.
% \jiaman{This is achieved by RefineNet, right?}\zhen{Yes, the wrist pose condition is achieved by RefineNet. But here I also want to stress that the finger pose should stay to $\hat{\bm{\theta}}$ during the contact phase.} \jiaman{I see. So basically here, you want to say we use the optimized grasp pose for all the timesteps in the contact phase? Then we should explicitly say this.}\zhen{Fixed}

\subsubsection{RefineNet}
\label{sec:refine_net}

% Specifically, we encode the wrist pose and contact phase information by modifying the input masked data conditions (introduced in Section 5.1). We expand the data dimension to include an additional term representing the relative wrist pose. For the frames in the contact phase, the optimized wrist pose $\hat{\bm{w}} \in \mathbb{R}^{T\times18}$ includes the relative position and 6D orientation of both wrists. For the frames in the pre-contact and post-contact phases, we pad this term with zeros.

The goal of RefineNet is to re-generate human-object interaction motions that align with the optimized grasp pose $\hat{\bm{g}}$. We adopt a conditional diffusion model based on CoarseNet, adding two conditions: \textit{wrist-object relative pose} to align the wrist with $\hat{\bm{w}}$ and \textit{object static} to keep the object stationary during non-contact phases. RefineNet’s input condition, $\bm{c}_{r} = \left\{\bm{W}, \bm{S}_{r}, \bm{G}, \bm{T}\right\}$, includes $\bm{G}$ and $\bm{T}$ from CoarseNet, with $\bm{W}$ and $\bm{S}_{r}$ detailed as follows. 


\vspace{-3mm}
\paragraph{Wrist-Object Relative Pose Condition.} 
The condition $\bm{W} \in \mathbb{R}^{T \times 18}$ represents the position and 6D rotation of both wrists in the object's frame, including only contact-phase wrist poses, with other entries set to zero.
% During training, it only includes the wrist pose for the contact phase, while all other entries are masked with zeros.
% During sampling, we set the entries corresponding to the contact phase with $\hat{\bm{w}}$, and all other entries remain zero.
% \jiaman{we can say contact phase?}\zhen{Here I’m referring to the training process. “contact phase” is the term specific to sampling, so I tried to avoid using it in this context. What do you think?} \jiaman{But there's no difference when you compute contact phases in training and sampling? Maybe we can just use contact phase. It's a bit confusing now.}\zhen{Fixed}

\vspace{-2mm}
\paragraph{Object Static Condition.}
In CoarseNet, the masked motion condition $\bm{S}$ was zero-padded except at the start, end, and waypoints. RefineNet uses contact phase information from CoarseNet to adjust $\bm{S}$ to $\bm{S_{r}}$, assigning a static object pose in pre-contact and post-contact frames.
% We apply the same approach during sampling.
% \jiaman{In training, you have GT data. GT data should already have static object poses in pre-contact and post-contact phases? Why do we need to adjust the object pose?}\zhen{Because in CoarseNet (\cref{sec:coarse_net}) training, $\bm{S}$ only contains the initial and end object pose. All other entries, including pre-contact and post-contact phases, are masked with zero. In RefineNet training, we remove the mask on the pre-contact and post-contact parts, so that the model knows how to deal with this static condition during sampling.}.\jiaman{Got it. But it seems training and sampling are the same? so we don't need to repeat this sentence. }\zhen{Fixed}

% Now, with $\bm{c} = \left\{\bm{S}, \bm{G}, \bm{T}, \bm{W}\right\}$, we further introduce a loss term to enforce the wrist-object relative pose condition.

\vspace{-2mm}
\paragraph{Wrist-Object Relative Pose Loss.}
We incorporate an additional loss to enforce the wrist-object relative pose condition. We sample 100 points $\bm{K}_{\text{rest}} \in \mathbb{R}^{100 \times 3}$ uniformly on the rest object’s surface and pre-calculate their positions in both wrists' coordinate frames, denoted as $\bm{K}_{\text{w}} \in \mathbb{R}^{2\times T \times 100 \times 3}$. 
At each time step, we compute the points' positions in the wrist frame, $\hat{\bm{K}}_{\text{w}}$, using the predicted orientations and positions of the object and wrist, $\bm{R}_{\text{o}}, \bm{T}_{\text{o}}, \bm{R}_{\text{w}}, \bm{T}_{\text{w}}$:
\begin{align}
\bm{K}_{\text{global}} &= \bm{R}_{\text{o}}\bm{K}_{\text{rest}} + \bm{T}_{\text{o}}, \\
\hat{\bm{K}}_{\text{w}} &= \bm{R}_{\text{w}}^{-1} \left( \bm{K}_{\text{global}} - \bm{T}_{\text{w}}\right).
\end{align}
The loss is computed as:
\begin{align}
\mathcal{L}_{\text{relative }}=
\sum_{t=1}^T
\bm{L}_{t}
\left\| 
\hat{\bm{K}}_{\text{w}, t}-\bm{K}_{\text{w}, t}
\right\|_1,
\end{align}
where $\bm{L}_{t}$ represents the contact labels, masking out the loss when the hand and object are not in contact.

% In addition, we adopt foot contact loss~\cite{tevet2022human} to mitigate foot sliding.



\vspace{-2mm}
\paragraph{Post-Processing.}
RefineNet generates human-object motions that align closely with input conditions, though minor adjustments may be required. To refine alignment, we apply post-processing: replacing object poses in pre- and post-contact phases with static poses and recalculating wrist trajectories during the contact phase based on object motion and the optimized wrist pose $\hat{\bm{w}}$. Additional implementation details are in the supplemental materials.







\subsubsection{FingerNet}
\label{sec:low_level_finger}
In this stage, we generate finger motions for the pre-contact and post-contact phases to create natural transitions between the rest and grasp poses, forming a complete interaction sequence along with predictions in prior stages.
Given the start and end finger poses and the wrist trajectory in between, we employ a conditional diffusion model to predict the finger motions denoted as $\bm{F} \in \mathbb{R}^{T\times D'}$, representing local 6D rotations~\cite{zhou2019continuity}.

% Similar to ManipNet~\cite{zhang2021manipnet}, we only process the finger motions for the right hand. 
% To handle the left hand, we mirror the entire scene, process the mirrored left hand as if it were the right hand, and then mirror the output back.
% This approach helps counteract the limited scale of the dataset~\cite{GRAB:2020}.
% At training time, both hand data are given so there is no bias on handedness.

% \paragraph{Motion Representation} 
% Finger motion $\boldsymbol{F} \in \mathbb{R}^{T\times D'}$ contains local joint rotations in 6D representation. 

% \paragraph{Hand-Object Spatial Representation.} 


 
% \paragraph{Input Condition Representation.} 
The model is conditioned on $\bm{c}_{f} = \left\{\bm{P}, \bm{F}_{s}, \bm{F}_{e}\right\}$, where $\bm{P}$ captures spatial hand-object relationships and $\bm{F}_{s}, \bm{F}_{e}$ are the start and end finger poses. Following prior work~\cite{zhang2021manipnet,taheri2023grip}, $\bm{P} \in \mathbb{R}^{T \times 100}$ is computed by sampling 100 palm-side mesh vertices and measuring their closest distances to the object, with the finger pose set to mean pose when calculating $\bm{P}$. 
% \jiaman{Is the data representation P the same as prior work? We probably need to clearly say what is prior work already have and what is new thing we proposed. So is P the same, but conditional diffusion model with P is new?} \zhen{Yes P is the same, but conditional diffusion model with P is new. Fixed.}


% \paragraph{Mirroring Strategy.} 
% The FingerNet is trained to process only the right hand. To generate finger motions for the left hand, we apply a mirroring strategy similar to~\cite{zhang2021manipnet}.
% Specifically, we first mirror the left hand and the corresponding object mesh to simulate a right-hand scenario. We then process this mirrored data to generate the finger motions. Finally, we mirror the resulting motions back to obtain the finger movements of the left hand. 

% To obtain the training data, we extracted all hand-reaching and hand-leaving motions from the hand-object interaction dataset. 
% The model is then optimized with the "simple" objective \cite{ho2020denoising} as described in Equation \ref{equ:simple}.

% During inference, we utilize the wrist trajectory generated by the interaction module, along with the resting finger pose and the precise finger pose $\hat{\boldsymbol{\theta}}$, as conditions for generating finger movements.









\subsection{Navigation Module}
\label{sec:low_level_navigation}
To generate long sequences of interactions with multiple objects, the human needs to navigate through the scene based on waypoints.
We adopt a conditional diffusion model to generate human navigation motion $\bm{H} \in \mathbb{R}^{T \times D}$. The input conditions consist of the initial human pose, 2D waypoints, and waypoint orientations (expressed as normalized direction vectors) to maintain consistent facing directions. 
Smooth transitions between the interaction and navigation modules are achieved by using the final pose of one module as the starting pose for the next.

% This module utilizes a conditional diffusion model to predict human motion, represented as $\bm{H} \in \mathbb{R}^{T \times D}$, and feet-ground contact labels, denoted as $\bm{L_f} \in \mathbb{R}^{T \times 2}$. The model conditions on the initial human pose, a set of sparse waypoints, and textual inputs. To generate motions with consistent facing directions between every two waypoints, the condition $\bm{c}$ integrates both the position and orientation of the waypoints, with orientation expressed as a normalized direction vector pointing from the current to the next waypoint. To ensure smooth transitions between the interaction and navigation modules, the final pose from one module is used as the starting pose for the subsequent module.


% The model architecture is the same as the interaction module, but trained with only the "simple" object \cite{ho2020denoising}:
% \begin{align}
%     \mathcal{L} &= \mathbb{E}_{x_0, n}\left\|\boldsymbol{F}_{\text{Navi}}\left(x_n, n, c \right)-x_0\right\|_1.
% \end{align}

% \subsection{Transition Between Interaction and Navigation}

% To achieve natural motion generation, we aim for sufficiently smooth transitions between the interaction module and navigation module. 
% Though both modules are conditioned on the initial human pose, simply adopting the last pose of one module as the initial pose for the next module does not suffice for two reasons.
% Firstly, this approach only ensures that poses are proximate, but it fails to ensure continuity in velocity. Secondly, since the two diffusion models are trained on distinct datasets with differing data distributions, directly setting the output pose of one module as the initial pose for another module does not yield satisfactory results.

% Therefore, we adopt the following methods to achieve smooth transitions.

% \medskip\textit{Interaction to Navigation Transition. } 
% 1. cut static ?????
% The end pose of the interaction module can vary widely, including poses such as bending over or raising hands above the head. Directly setting these poses as the initial pose of the navigation module can result in unnatural movement, as the model hasn't encountered such motions during training. 
% To tackle this issue, we create a new initial pose by substituting the upper body pose with a rest pose while preserving the lower body pose, as shown in Figure XXX. We then utilize interpolation techniques to ensure a smooth upper body transition.
% We empirically observe that this approach significantly improves the generation of natural movement.

% \medskip\textit{Navigation to Interaction Transition. } 
% We found that the interaction module effectively handles initial poses where the human is close to the object, especially with walking or standing poses. Thus, we switch to the interaction module when the human is within a distance range of 0.3-0.5 meters from the object. 
% Additionally, we empirically observe that transitioning when both feet are on the ground helps maintain velocity continuity and results in more natural motion.


\section{Physics Tracker}
After obtaining a sequence of kinematic motions, we train a tracking policy using reinforcement learning to control a physically simulated character that imitates these motions. This approach helps eliminate artifacts (e.g., hand-object penetration and foot sliding) in the kinematic motions and ensures the physical plausibility of the generated results.

We use PPO~\cite{schulman2017proximal} as the backbone reinforcement learning algorithm,
and adopt an importance sampling strategy to facilitate the tracking of long motion sequences involving interactions with multiple objects.
Policies trained by our approach can accurately track the kinematics motions, enabling the character to grasp and move distinct objects along the generated trajectories.
Implementation and training details for the physics tracker are provided in the supplementary materials.
\section{Experiment}
\label{sec:experiment}
We conduct extensive evaluation to assess the overall performance of our system, as well as each individual component. We use quantitative metrics and user perception studies to evaluate our results compared to the baselines. We also conduct ablation studies to analyze the effects of each system module. 
We present the key results in this section and highly recommend readers refer to the supplementary video and documents for more detailed results and descriptions of the studies.
% We present the key results in this section. More results and detailed descriptions of the studies can be found in the supplementary video and documents.
% We evaluate each system module by comparing our proposed approach with baseline methods and conducting ablation studies. 
% We highly recommend that readers refer to the supplementary video for additional qualitative results.

\subsection{Evaluation of High-Level Planner}
\label{subsec:planner_eval}

% \karen{Geometric accuracy: Describe metrics. Describe samples. Refer to a table or figure. Conclusion.}

% \karen{Perception study: Describe metrics. Describe samples. Refer to a table or figure. Conclusion.}

The core component of the high-level planner is the scene map calculation. To assess its effectiveness, we quantitatively measure the geometric accuracy of the resulting plans and conduct a human perception study.
For evaluation, we compare our method to a baseline that uses LLMs to directly predict object positions and orientations. We design 25 human-level instructions covering common everyday tasks for assessment, using OpenAI's GPT-4o \cite{gpt4o} for both our method and the baseline.

\begin{figure}[t]
\begin{center}
    \includegraphics[width=\columnwidth]{figures/4_llm_results_compressed.pdf}
    \vspace{-5mm}
    \caption{Comparison between (a) the baseline and (b) our high-level planner. The baseline generates an incorrect monitor position and orientation, and places the chair in a position that intersects with the table. The human perception study results (bottom) show that the majority of participants preferred our results.
    %\karen{Should show the input text commands too. Also need caption on the perception study.}\zhen{Added.}
    }
    \label{fig:llm_result}
   \vspace{-8mm}
\end{center}
\end{figure}



\begin{table*}[ht]
\centering
\begin{adjustbox}{max width=\textwidth}
\begin{tabular}{lcccccccccccccccccc}
\toprule
\multirow{2}{*}{Method} & \multicolumn{1}{c}{Condition} & \multicolumn{2}{c}{Human Motion} & \multicolumn{5}{c}{Interaction} & \multicolumn{3}{c}{GT Difference} \\
\cmidrule(r){2-2} \cmidrule(r){3-4} \cmidrule(r){5-9} \cmidrule(r){10-12}
& $T_{xy} \downarrow$ & $H_{\text{feet}} \downarrow$ & FS $\downarrow$ & $\text{C}_{\text{prec}} \uparrow$ & $\text{C}_{\text{rec}} \uparrow$ & $\text{C}_{\text{F1}} \uparrow$ & CC $\uparrow$ & IV $\downarrow$ & $\text{MPJPE} \downarrow$ & $T_{\text{obj}} \downarrow$ & $O_{\text{obj}} \downarrow$ \\
\midrule
% IMoS \cite{ghosh2023imos} 
% & 112.32  & -90.32 & 1.46 & 0.26 & 0.07 & 0.09 & 1.25\% & 18.25 & 32.86 & 130.10 & 2.36 \\
CNet+GRIP \cite{taheri2023grip} 
& 3.58 & \textbf{2.50} & 0.45 & 0.90 & 0.64 & 0.70 & 3.99\% & 49.00 & 14.33 & 11.55 & 0.93 \\
\midrule
CNet 
& 3.26 & 2.68 & 0.43 & 0.91 & 0.81 & 0.84 & 3.54\% & 48.56 & \textbf{14.32} & \textbf{10.56} & 1.06 \\
C+RNet 
& \textbf{2.81} & 2.89 & \textbf{0.33} & \textbf{0.91} & \textbf{0.95} & \textbf{0.92} & 3.84\% & 24.78 & 15.96 & 12.57 & \textbf{0.73} \\
C+R+FNet (ours)
& \textbf{2.81} & 2.89 & \textbf{0.33} & \textbf{0.91} & \textbf{0.95} & \textbf{0.92} & \textbf{5.53\%} & \textbf{19.06} & 15.96 & 12.57 & \textbf{0.73}\\
\bottomrule
\end{tabular}
\end{adjustbox}
\caption{Interation synthesis on the FullBodyManipulation dataset \cite{li2023object}. Our full model achieves superior interaction quality compared to all others. }
\label{tab:interaction}
\end{table*}


\begin{figure*}[t]
    \centering
     % \captionsetup{aboveskip=0pt, belowskip=0pt}
\vspace{-2mm}
\includegraphics[width=\textwidth]{figures/5_qualitative_results_v4_compressed.pdf}
\vspace{-8mm}
    \caption{Qualitative comparisons of each method. All baselines and ablations fail to produce precise hand and finger motions, while our method generates natural hand-object interactions.} 
    \label{fig:interaction_qualitative}
\vspace{-4mm}
\end{figure*}

\begin{figure}[t]
    \centering
\includegraphics[width=0.48\textwidth]{figures/6_interaction_human_study_v2.pdf}
\vspace{-5mm}
    \caption{Comparison of our full model and baseline models through human perception studies.}
    \label{fig:interaction_human_study}
\vspace{-6mm}
\end{figure}




\setlength{\columnsep}{10pt}
\begin{wraptable}{r}{0.45\columnwidth}
\centering
% \begin{table}[]
% \centering
\vspace{-10pt}
\caption{Comparison of geometric accuracy between our method and the baseline.}
\begin{adjustbox}{max width=0.45\columnwidth}
\begin{tabular}{lcc}
\toprule
& PE$_\text{p} \downarrow$ & PE$_\text{o} \downarrow$ \\
\midrule
Baseline & 21.9\% & 12.5\% \\
Ours & \textbf{3.1\%} & \textbf{1.6\%} \\
\bottomrule
\end{tabular}
\end{adjustbox}
\label{tab:llm_geometry}
\vspace{-10pt}

\end{wraptable}
% \end{table}
For geometric accuracy, we measure the proportion of objects with positional (PE$_\text{p}$) or orientation errors (PE$_\text{o}$) in the scene. 
Positional errors include objects placed at incorrect heights or with penetration into other objects, while orientation errors refer to objects being directed in the wrong direction.
The results, shown in \cref{tab:llm_geometry}, demonstrate that our method outperforms the baseline by a large margin. A qualitative comparison is presented in the upper part of \cref{fig:llm_result}, where the baseline generates incorrect object positions and orientations, while our method produces more reasonable layouts.



% \paragraph{Results.} 
We also conducted a human perception study, with results presented in the lower part of \cref{fig:llm_result}. Each pair was reviewed by 20 Amazon Mechanical Turk (AMT) workers. They were asked to choose the scene that was more reasonable and consistent with the given instruction. The majority of participants preferred our method.

\subsection{Evaluation of Low-Level Motion Generator}
\label{subsec:low_level_eval}
% \karen{The low-level motion generator is evaluated by 11 metrics, including two newly introduced metrics for finger-object interaction and other frequently used metrics in the area of human motion synthesis. Additionally, we conduct a perception study to validate the quality of synthesized motion that might be difficult to quantify by numerical metrics. We report the results in comparison with the baseline and ablations to analyze the overall impact of our system, as well as the contributions of each individual component.}


\paragraph{Datasets.}
CoarseNet and RefineNet are trained on the FullBodyManipulation dataset~\cite{li2023object}, which includes 10 hours of human-object interaction with 15 objects but lacks finger motion.
FingerNet uses the GRAB dataset~\cite{GRAB:2020}, containing full-body and finger motion for 10 subjects with 51 objects.
The navigation module is trained on HumanML3D~\cite{guo2022generating}, featuring 28 hours of diverse motions with language descriptions.
For all datasets, we follow standard data partitioning from prior work~\cite{li2023object, GRAB:2020, guo2022generating}.

\vspace{-4mm}
\paragraph{Evaluation Metrics.}
Following prior work~\cite{li2023controllable}, we evaluate our approach from various aspects.
All distance errors are measured in centimeters. 
For condition matching, we report the waypoint following errors ($T_{xy}$).
For human motion quality, we report the foot sliding score (FS) and foot height ($H_{\text{feet}}$).
% We also report FID, which calculates the discrepancy between the generated and ground truth motions. Additionally, R-precision (R\textsubscript{prec}) measures how well the generated motion matches the input text by checking if the correct motion is ranked among the top-3 results for a given text description.
For interaction quality, we use contact precision (C\textsubscript{prec}), recall (C\textsubscript{rec}), and F1 score (C\textsubscript{F1}).
To better assess finger-object interaction, we also report contact coverage (CC) -- the percentage of hand points within -2mm to +2mm of the object surface -- and intersection volume (IV), the overlap volume of hand and object meshes in $\text{cm}^3$, following~\cite{grady2021contactopt}.
For ground truth difference, we report the mean per-joint position error (MPJPE), object position error ($T_{\text{obj}}$), and the object orientation error ($O_{\text{obj}}$).
We also conduct a human perception study to measure the motion quality.

\vspace{-4mm}
\paragraph{Baselines.}
For generating full-body human motion, finger motion, and object motion from text, prior work~\cite{ghosh2022imos,li2024task} trained models on GRAB~\cite{GRAB:2020} to predict motions for small-object manipulation. However, GRAB contains only small-object manipulation data without locomotion, so models trained on it cannot generalize to scenarios involving large objects that require full-body coordination and locomotion. Moreover, the approach from~\cite{ghosh2022imos,li2024task} is not applicable to the FullBodyManipulation dataset~\cite{li2023object} used in our work, as this dataset lacks finger motions, making it infeasible to learn a model via direct supervision. 
Due to these differences, a directly comparable baseline does not exist.
Thus, we establish our baseline by combining two complementary methods: CHOIS~\cite{li2023controllable} (referred to as CoarseNet in our framework), which generates object motion and full-body motion from text, and GRIP~\cite{taheri2023grip}, which refines arm movements and generates finger motions based on CHOIS's output.

In addition, to evaluate the impact of each component, we compare our full system against two ablations: (1) CoarseNet alone (CNet) and (2) CoarseNet + RefineNet (C+RNet), which represents the full system without finger motion. 
The full system is denoted as C+R+FNet.
% Our full system, which can be seen as (C+RNet) combined with FingerNet, is denoted as C+R+FNet.
This evaluation focuses on the motion synthesis quality before physics tracking. We present the evaluation for the physics tracker in \cref{subsec:phys_eval}.
%We evaluate all methods on the FullBodyManipulation dataset~\cite{li2023object}. 


\vspace{-4mm}
\paragraph{Results.}
The qualitative comparison in \cref{fig:interaction_qualitative} clearly demonstrates the superiority of our method. Both ablated variants show a lack of finger detail, and GRIP~\cite{taheri2023grip} fails to generate accurate hand motion, resulting in significant artifacts. In contrast, our method produces precise finger-object interactions.
We also observe the impact of RefineNet: compared to CNet, C+RNet maintains more consistent and reasonable relative poses between the hand and object during the interaction, showing RefineNet’s role in improving pose consistency and reducing artifacts.

The quantitative results are presented in \cref{tab:interaction}.
Since our method focuses on finger-object interaction, the five metrics under the \textit{Interaction} column are particularly important, and we achieve the best performance across all of them.

We further conduct a human perception study to complement our evaluation. We generate 15 sequences using each method and compose 60 pairs, with each pair consisting of our results and the results of a baseline. Each pair was reviewed by 20 Amazon Mechanical Turk (AMT) workers. As shown in \cref{fig:interaction_human_study}, our method is consistently preferred over the baselines by a large margin. Notably, our method achieves even higher preference rates when compared to the real-world data from the FullBodyManipulation~\cite{li2023object} dataset, as we generate detailed finger motions that are absent in the original dataset.

\subsection{Evaluation of Physics Tracker}
\label{subsec:phys_eval}

The physics tracker can accurately track kinematic motion while ensuring the physical plausibility of the motion. Through physics simulation, artifacts such as foot floating, foot sliding, and human-object penetration are eliminated, as shown in \cref{fig:physics_qualitative_results}. 

\setlength{\columnsep}{10pt}
\begin{wraptable}{r}{0.4\columnwidth}
\centering
% \begin{table}[]
% \centering
\vspace{-1pt}
\caption{Tracking error between the kinematic motion and physics-based motion.}
\begin{adjustbox}{max width=0.35\columnwidth}
\begin{tabular}{lcc}
\toprule
& $E_h$ & $E_o$ \\
\midrule
Ours & 5.45 & 4.67 \\
\bottomrule
\end{tabular}
\end{adjustbox}

\label{tab:physics}
\vspace{-6pt}
\end{wraptable}
% \end{table}
Additionally, we evaluate the tracking error of human joints ($E_h$) and objects ($E_o$), as shown in \cref{tab:physics}. The tracking error is defined as the positional error (cm) between the kinematic motion and the tracked motion. 
The results were evaluated across 20 interaction sequences, each with an average duration of 30 seconds and 2 objects of varying shapes and types.
All sequences were successfully tracked, with relatively small errors, demonstrating high tracking accuracy.



\begin{figure}[t]
    \centering
     % \captionsetup{aboveskip=0pt, belowskip=0pt}
\includegraphics[width=0.48\textwidth]{figures/7_physics_qualitative_results_v3_compressed.pdf}
\vspace{-7mm}
    \caption{Qualitative comparison of kinematic motion and physics-based motion. The physics tracker corrects artifacts such as foot floating and penetration.} 
    \label{fig:physics_qualitative_results}
    % \vspace{-8mm}
\end{figure}


\section{Conclusions}
In this work, we introduced a system that simulates realistic human-object manipulation in a contextual environment based on human-level language instructions.
The framework consists of a high-level LLM planner that infers a target scene layout and execution plan from the given instructions, a low-level motion generator that generates synchronized, full-body, fingers, and object motion, and a physics tracker that imitates the generated kinematic motions and produces physically plausible interaction motions.
Our system effectively handles a variety of tasks, showcasing its potential for various real-world applications.

\paragraph{Future Work.}
% Despite its strengths, the system has several limitations. 
% While our approach successfully leverages LLMs to interpret human-level instructions, there is room to improve plan generation for more complex environments.
Our system currently uses text or image inputs to represent the contextual environment, and could potentially be enhanced with more structured representations like voxel grids as used in \cite{jiang2024autonomous}. 
% For grasp pose generation, we build upon~\cite{wang2023dexgraspnet}, though alternative approaches could be explored to further improve grasp success rates.
% \zhen{Include failure cases in the video.}
% \paragraph{Future Work.}
Building on this idea, a promising direction is to create an agent with egocentric vision perception as input to LLM/VLMs for high-level motion planning. 
This could enable the agent to exhibit more natural behaviors, such as purposeful head orientation for observation and better body positioning to avoid collisions.

% A promising future direction is to create an agent with egocentric vision perception and use it as an input to the LLM/VLMs for high-level motion planning. We expect that more humanlike behaviors will emerge as a result of realistic modeling of human perception. For example, the agent will orient its head more purposefully to observe the world and will position its whole body more effectively to avoid collisions. 

% \paragraph{Future Work.}
% First, we aim to improve the quality and variation of generated motions for more realistic interactions. We also plan to enhance the system's ability to handle more complex environments and diverse interactions, such as sitting on a chair or lying on a sofa \jiaman{This is contradictory to what we said in the introduction. sitting on a chair is not complex, it's solved.}.\zhen{But our current system can't do such tasks, I want to say that we can build a unified system that can handle both static and dynamic objects,}
\paragraph{Acknowledgement.}
We thank Wenyang Zhou for data processing. This work is in part supported by the Wu Tsai Human Performance Alliance at Stanford University and the Stanford Institute for
Human-Centered AI (HAI). 

% Bibliography
{
    \small
    \bibliographystyle{ieeenat_fullname}
    \bibliography{sample-bibliography}
}


% WARNING: do not forget to delete the supplementary pages from your submission 
\clearpage
\newpage
% CVPR 2025 Paper Template; see https://github.com/cvpr-org/author-kit

% \documentclass[10pt,twocolumn,letterpaper]{article}

% %%%%%%%%% PAPER TYPE  - PLEASE UPDATE FOR FINAL VERSION
% \usepackage{iccv}              % To produce the CAMERA-READY version
% \usepackage[review]{iccv}      % To produce the REVIEW version
% \usepackage[pagenumbers]{cvpr} % To force page numbers, e.g. for an arXiv version

% Import additional packages in the preamble file, before hyperref
% %
% --- inline annotations
%
\newcommand{\red}[1]{{\color{red}#1}}
\newcommand{\todo}[1]{{\color{red}#1}}
\newcommand{\TODO}[1]{\textbf{\color{red}[TODO: #1]}}
% --- disable by uncommenting  
% \renewcommand{\TODO}[1]{}
% \renewcommand{\todo}[1]{#1}




% Include other packages here, before hyperref.

% It is strongly recommended to use hyperref, especially for the review version.
% hyperref with option pagebackref eases the reviewers' job.
% Please disable hyperref *only* if you encounter grave issues, e.g. with the
% file validation for the camera-ready version.
%
% If you comment hyperref and then uncomment it, you should delete
% ReviewTempalte.aux before re-running LaTeX.
% (Or just hit 'q' on the first LaTeX run, let it finish, and you
%  should be clear).
% \definecolor{cvprblue}{rgb}{0.21,0.49,0.74}
% \usepackage[pagebackref,breaklinks,colorlinks,allcolors=cvprblue]{hyperref}
% \AtBeginDocument{
%   \addtocontents{toc}{\protect\hypersetup{linkcolor=black}}
% }
% \AtEndDocument{
%   \addtocontents{toc}{\protect\hypersetup{linkcolor=cvprblue}}
% }
% \usepackage{verbatim}
% \usepackage{comment}
% \usepackage{wrapfig}

% \usepackage{graphicx}
% \usepackage{amsmath}
% \usepackage{amssymb}
% \usepackage{booktabs}
% % \usepackage{caption}
% \usepackage{bm}
% \usepackage{adjustbox}
% \usepackage{multirow}

% \usepackage{url}            % simple URL typesetting
% \usepackage{amsfonts}       % blackboard math symbols
% \usepackage{nicefrac}       % compact symbols for 1/2, etc.
% \usepackage{microtype}      % microtypography
% \usepackage{xcolor}         % colors

% \usepackage{array}
% \usepackage{xspace}

% \usepackage{float}
% Support for easy cross-referencing
% \usepackage[capitalize]{cleveref}
% \crefname{section}{Sec.}{Secs.}
% \Crefname{section}{Section}{Sections}
% \Crefname{table}{Table}{Tables}
% \crefname{table}{Tab.}{Tabs.}


% \usepackage[ruled]{algorithm2e} % For algorithms



% % require xspace, array

\newcommand{\stderr}[1]{\scriptsize $\pm #1$}

%% layout
\newcolumntype{L}[1]{>{\raggedright\let\newline\\\arraybackslash\hspace{0pt}}m{#1}}
\newcolumntype{C}[1]{>{\centering\let\newline\\\arraybackslash\hspace{0pt}}m{#1}}
\newcolumntype{R}[1]{>{\raggedleft\let\newline\\\arraybackslash\hspace{0pt}}m{#1}}

\newcommand{\xpar}[1]{\noindent\textbf{#1}\ \ }
\newcommand{\vpar}[1]{\vspace{3mm}\noindent\textbf{#1}\ \ }

%% notations
\newcommand{\sect}[1]{Section~\ref{#1}}
\newcommand{\sectapp}[1]{Appendix~\ref{#1}}
\newcommand{\ssect}[1]{\S~\ref{#1}}
\newcommand{\eqn}[1]{Equation~\ref{#1}}
\newcommand{\fig}[1]{Figure~\ref{#1}}
\newcommand{\tbl}[1]{Table~\ref{#1}}
\newcommand{\alg}[1]{Algorithm.~\ref{#1}}

\newcommand{\degree}{\ensuremath{^\circ}\xspace}
\newcommand{\ignore}[1]{}
\newcommand{\norm}[1]{\lVert#1\rVert}
\newcommand{\fcseven}{$\mbox{fc}_7$}
% \DeclarePairedDelimiter\ceil{\lceil}{\rceil}
% \DeclarePairedDelimiter\floor{\lfloor}{\rfloor}

\renewcommand*{\thefootnote}{\fnsymbol{footnote}}

%\DeclareMathOperator*{\argmin}{arg\,min}
%\DeclareMathOperator*{\argmax}{arg\,max}

\def\naive{na\"{\i}ve\xspace}
\def\Naive{Na\"{\i}ve\xspace}

\makeatletter
\DeclareRobustCommand\onedot{\futurelet\@let@token\@onedot}
\def\@onedot{\ifx\@let@token.\else.\null\fi\xspace}

\def\eg{e.g\onedot} \def\Eg{E.g\onedot}
\def\ie{i.e\onedot} \def\Ie{I.e\onedot}
\def\cf{\emph{c.f}\onedot} \def\Cf{\emph{C.f}\onedot}
\def\etc{etc\onedot}
\def\vs{\emph{vs}\onedot}
\def\wrt{w.r.t\onedot}
\def\dof{d.o.f\onedot}
\def\etal{et al\onedot}
\def\aka{a.k.a\onedot}
\makeatother

%% comments
\definecolor{MyDarkBlue}{rgb}{0,0.08,1}
\definecolor{MyDarkGreen}{rgb}{0.02,0.6,0.02}
\definecolor{MyDarkRed}{rgb}{0.8,0.02,0.02}
\definecolor{MyDarkOrange}{rgb}{0.40,0.2,0.02}
\definecolor{MyPurple}{RGB}{111,0,255}
\definecolor{MyRed}{rgb}{1.0,0.0,0.0}
\definecolor{MyGold}{rgb}{0.75,0.6,0.12}
\definecolor{MyDarkgray}{rgb}{0.66, 0.66, 0.66}

\definecolor{turquoise}{cmyk}{0.65,0,0.1,0.3}

\newcommand{\karen}[1]{\textcolor{red}{[Karen: {#1}]}}
\newcommand{\jiaman}[1]{\textcolor{orange}{[Jiaman: {#1}]}}
\newcommand{\zhen}[1]{\textcolor{blue}{[Zhen: {#1}]}}

\renewcommand{\algorithmcfname}{ALGORITHM}

% %%%%%%%%% PAPER ID  - PLEASE UPDATE
% \def\paperID{6622} % *** Enter the CVPR Paper ID here
% \def\confName{ICCV}
% \def\confYear{2025}


% % Document starts

% % Title portion
% \title{Human-Object Interaction from Human-Level Instructions}

% \author{First Author\\
% Institution1\\
% Institution1 address\\
% {\tt\small firstauthor@i1.org}
% % For a paper whose authors are all at the same institution,
% % omit the following lines up until the closing ``}''.
% % Additional authors and addresses can be added with ``\and'',
% % just like the second author.
% % To save space, use either the email address or home page, not both
% \and
% Second Author\\
% Institution2\\
% First line of institution2 address\\
% {\tt\small secondauthor@i2.org}
% }

% \begin{document}

% \twocolumn[{%
% \renewcommand\twocolumn[1][]{#1}%
% \maketitle
% \begin{center}
%     \centering
% \includegraphics[width=\linewidth]{figures/0_teaser.pdf}
% % \vspace{-2em}
% \captionof{figure}{\zhen{need revision}}
% \label{fig:teaser}
% \end{center}
% }]

% \begin{document}


% WARNING: do not forget to delete the supplementary pages from your submission 


% \clearpage
\renewcommand{\thefigure}{S\arabic{figure}}
\renewcommand{\thetable}{S\arabic{table}}
\def\theequation{S\arabic{equation}}
\setcounter{page}{1}
\setcounter{section}{0}
\setcounter{equation}{0}
\setcounter{table}{0}
\setcounter{figure}{0}
\maketitlesupplementary

\renewcommand{\thesection}{\Alph{section}}
\renewcommand{\thesubsection}{\thesection.\arabic{subsection}}
% \clearpage
% \newpage
\setcounter{tocdepth}{2}
% \tableofcontents



\section{Details of High-Level Planner}



\subsection{Pseudo-code for Object Positions Calculation}

In \cref{alg:compute_position}, we provide the pseudo-code for the algorithm that calculates object positions in the scene graph, as described in~\cref{subsec:spatial_relationship}.

\begin{algorithm}
\SetAlgoLined
\DontPrintSemicolon
\SetKwInOut{Input}{Input}
\SetKwInOut{Output}{Output}

\Input{Scene graph $G(V, E)$ with vertices $V$ and edges $E$, initial positions $L$, static nodes ($V_{S}$), the nodes to be moved ($V_{M}$).}
\Output{Updated positions $L$ for all nodes.}

\SetKwFunction{FSub}{ComputePositions}
\SetKwFunction{FUpdate}{Update}
\SetKwFunction{FOffset}{ComputeOffset}

\SetKwProg{Fn}{Procedure}{:}{\KwRet}
\Fn{\FSub{$G$, $L$}}{
    % Let $E_{0}$ be the set of edges in $G_0$\;
    Let the set of processed nodes $V_{P} \gets V_{S}$\;
    
    \While{$V_{M}$ is not empty}{
        Let \( V' \subseteq V_M \) be the nodes whose predecessors' positions are known\;
        \ForEach{$v \in V'$}{
            % $L_{\text{P}}(u) \gets \{L(u) \mid (u, v) \in E\}$\;
            $L(v) \gets$ \FUpdate{$v$, $G$, $L$}\;
            $V_{P} \gets V_{P} \cup \{v\}$\;
            $V_M \gets V_M \setminus \{v\}$
            % Remove $v$ and all edges $(u, v)$from $G$\;
        }
    }
    \KwRet $L$\;
}

\Fn{\FUpdate{$v$, $G$, $L$}}{
    Initialize an empty list of positions $P$\;
    \ForEach{predecessor $u$ of $v$ in $G$}{
        $O \gets$ \FOffset{$(u, v)$, $L(u)$}\;
        Append $L(u) + O$ to $P$\;
    }
    \KwRet $\text{Average}(P)$\;
}

\Fn{\FOffset{$(u, v)$, $L(u)$}}{
    \If{edge $(u, v)$ indicates ``on''}{
        Set $O_{\text{height}}$ to align with the top surface of $u$\;
        Sample $O_{\text{horizontal}}$ within the horizontal bounds of $u$\;
        \KwRet $[O_{\text{horizontal}}, O_{\text{height}}]$\;
    }
    \ElseIf{edge $(u, v)$ indicates ``adjacent''}{
        Use the direction and distance provided by the LLMs to compute $O$\;
        \KwRet $O$\;
    }
}
% $L \gets$ \FSub{$G$, $L$, $V_{S}$, $V_{M}$}\;

\caption{Object Positions Calculation}
\label{alg:compute_position}
\end{algorithm}

\subsection{Prompts for High-Level Planner}
We provide the prompts used for generating the scene map and execution plan in~\cref{fig:llm_spatial_text}. 
An example of the input and the corresponding output of the LLM planner is shown in~\cref{fig:llm_our_example}.
We also present the prompt used by the baseline (as described in~\cref{subsec:planner_eval})
% \cref{subsec:planner_eval}
in~\cref{fig:llm_baseline}.


\begin{figure*}[h]
    \centering
     % \captionsetup{aboveskip=0pt, belowskip=0pt}
\includegraphics[width=\textwidth]{figures/supp_our_prompt.pdf}
    \caption{Prompts used by the LLM planner to generate the scene map and execution plan.} 
    \label{fig:llm_spatial_text}
    \vspace{-2mm}
\end{figure*}


\begin{figure*}[h]
    \centering
     % \captionsetup{aboveskip=0pt, belowskip=0pt}
\includegraphics[width=\textwidth]{figures/supp_input_output.pdf}
    \caption{Example of the input and output for the LLM planner.} 
    \label{fig:llm_our_example}
    \vspace{-2mm}
\end{figure*}


\begin{figure*}[h]
    \centering
     % \captionsetup{aboveskip=0pt, belowskip=0pt}
\includegraphics[width=\textwidth]{figures/supp_baseline_prompt.pdf}
    \caption{Prompts used by the baseline, as described in~\cref{subsec:planner_eval}.} 
    \vspace{-2mm}
    \label{fig:llm_baseline}
\end{figure*}


\section{Details of Low-Level Motion Generator}
\subsection{Diffusion Model Architecture}
\label{supp:diff_arch}

Our system includes four diffusion models: CoarseNet, RefineNet, FingerNet, and the navigation module, all of which adopt a similar transformer-based architecture. Here, we illustrate the architecture of RefineNet in \cref{fig:supp_diffusion}. The model is provided with a motion sequence $\bm{x}_n$, the condition $\bm{c}_r$, and the noise step $n$, and predicts the clean motion $\hat{\bm{x}}_{0}$.
For the other networks, we simply substitute the condition with the corresponding one.

\begin{figure}[!t]
    \centering
     % \captionsetup{aboveskip=0pt, belowskip=0pt}
\includegraphics[width=0.48\textwidth]{figures/supp_diffusion.pdf}
    \caption{Architecture of RefineNet. The model is conditioned on $\bm{c}_{r} = \left\{\bm{W}, \bm{S}_{r}, \bm{G}, \bm{T}\right\}$, where $\bm{W}$ represents the wrist-object relative pose, $\bm{S}_{r}$ is the masked motion representation containing the static pose and waypoints, $\bm{G}$ denotes the object geometry, and $\bm{T}$ is the textual input.}
    \label{fig:supp_diffusion}
    
\end{figure}


\subsection{Implementation Details of Grasp Pose Generation}
As introduced in~\cref{subsubsec:grasp_pose}, we use a state-of-the-art grasp pose generation
method ~\cite{wang2023dexgraspnet} to generate the grasp pose.
Although originally designed for robotic hand grasping, this method relies solely on kinematic chain information, making it fully compatible with the SMPL-X~\cite{smplx} hand model. Starting with an object mesh and an initial hand pose, the grasp generation process is formulated as an optimization problem that minimizes a differentiable force closure energy term along with several additional terms that encourage the fingers to conform closely to the object surface while maintaining a natural hand shape. For motions involving both hands, such as carrying a box, we omit the force closure term.

\subsection{Post-processing in RefineNet}
\label{subsec:post_process}
As introduced in~\cref{sec:refine_net}, we apply post-processing to RefineNet's output to better align it with the input condition.
First, we replace the object poses in the pre-contact and post-contact phases with the corresponding static poses. To address the abrupt transitions at the phase boundaries, we identify the difference between the predicted object pose and the static object pose at the boundary frame denoted as $\Delta \bm{d} \in \mathbb{R}^6$ (suppose this is the boundary between pre-contact phase and contact phase), which encapsulates both the positional and rotational differences in axis-angle format. To smooth these transitions, we employ an interpolation strategy to the object motion in the contact phase. The interpolation function is defined by $\hat{\bm{O}}_{t} = \bm{O}_{t} + \alpha_t \Delta \bm{d}$, where $\alpha_t$ is a time-dependent scaling factor that decreases from 1 to 0 across the transition period, $\bm{O}_t$ denotes the predicted object pose. 

Next, we use the rectified object motion and the optimized wrist pose $\hat{\bm{w}}=(\bm{\hat{R}}, \bm{\hat{T}})$ to compute the wrist trajectories during the contact phase. 
The rotation is computed as $\bm{R}_o\bm{\hat{R}}$, and the position is computed as $\bm{R}_o \bm{\hat{T}} + \bm{T}_o$, where $\bm{R}_o$ and $\bm{T}_o$ represent the object's rotation and position, respectively.
The computed wrist trajectories may slightly deviate from the predicted wrist trajectories in the contact phase, leading to unsmooth transitions near the phase boundaries. We apply the same interpolation strategy to the wrist poses in the pre-contact and post-contact phases to produce smooth wrist motions. Finally, we employ an Inverse Kinematics (IK) algorithm to solve for the updated human pose constrained by the wrist poses. 


\subsection{Implementation Details of FingerNet}
FingerNet generates finger motions for the pre-contact and post-contact phases to create natural transitions between rest and grasp poses. It is trained on the GRAB dataset~\cite{GRAB:2020}. Since FingerNet focuses solely on the hand approaching and releasing motions, we only use the corresponding portions of the dataset. Specifically, we extract one second of motion before the grasp begins and after it ends as training data, discarding the rest.


We adopt a mirroring strategy similar to~\cite{zhang2021manipnet} to handle both hands simultaneously. During training, left-hand data are mirrored and combined with the right-hand data so that the model only needs to predict motions for the right hand. At inference time, we mirror the left-hand input, predict the corresponding motions using the model, and then mirror the output back to obtain the left-hand motions.

\subsection{Smooth Transition between Navigation and Interaction}
The navigation and interaction modules alternate to generate a long-horizon interaction sequence. Specifically, when the human is far from the object, the navigation module is activated to guide the human towards it. Once the human is within a threshold distance (set at 1 meter), the system switches to the interaction module to begin object manipulation. After the interaction is completed and the object is put down, control reverts to the navigation module to move towards the next object.

Since both modules use the initial human pose as an input condition, we ensure a smooth transition by feeding the last pose from the previous module into the next.
We apply the same smoothing technique described in \cref{subsec:post_process} at the switching boundary to further enhance continuity.


\section{Implementation Details of Physics Tracker}
We use IsaacGym~\cite{makoviychuk2021isaac} as the physics engine to perform simulation.
The simulated character is skeleton-driven where each joint is modeled as a motor controlled through a PD controller.
The control policy works at 30Hz, which is the same with generated kinematics motions, while the simulation runs at 120Hz.
Our character has $62$ body links and $49$ controllable joints.
This results in a state space $\mathbf{s}_t \in \mathbb{R}^{(62+|\mathcal{O}|)\times 13}$ including the position, orientation and linear and angular velocities of the character's body links and the manipulating objects $\mathcal{O}$ depending on the given tracking motions.
To track the target motion,
additional observation $\mathbf{o}_t \in \mathbb{R}^{(62+|\mathcal{O}|)\times7}$ including the target position and orientation of each body link and object is also introduced as the input to the control policy.
The action space is $\mathbf{a}_t \in \mathbb{R}^{49\times3}$ given each controllable joint has 3 degrees of freedom.
All our control policies are optimized using PPO~\cite{schulman2017proximal} as the backbone reinforcement learning algorithm with network parameters updated by Adam optimizers~\cite{kingma2014adam}.
The hyperparameters used for policy training are listed in Tab.~\ref{tab:rl_hyper}.

\subsection{Reward Design}
Since fingers take $30$ out of $49$ controllable joints and have a higher requirement in control precision for accurate object manipulation, we define a reward function containing three terms to evaluate the full-body pose and finger pose separately:
\begin{equation}
    r = 0.8 r_\text{body} + 0.2 r_\text{hand} + 0.05 r_\text{energy},
\end{equation}
where $r_\text{body}$ evaluates the full-body poses excluding fingers, $r_\text{hand}$ evaluates the hand poses related to the target manipulating object,
and $r_\text{energy}$ is an additional energy-related term penalizing end effectors' abrupt movement.
Note that in our implementation, at any time step,
at most one object is activated as the manipulation target, while multiple objects may exist in the scene and would be manipulated one by one.

\begin{table}[t]
    \centering
    \begin{tabular}{rl}
    \toprule
       \textbf{Body Part}  &  \quad\textbf{Value} \\
       \midrule
       root (pelvis) & $w_q = 1, w_p = 1$\\
       lower abdomen & $w_q = 0.2$\\
       upper abdomen & $w_q = 0.2$\\
       chest & $w_q = 0.2$\\
       neck & $w_q = 0.2$\\
       head & $w_q = 0.2$\\
       clavicles & $w_q = 0.1$\\
       upper arms & $w_q = 0.2$ \\
       lower arms & $w_q = 0.2$\\
       wrists & $w_q = 0.3, w_p = 0.3$\\
       thighs & $w_q = 0.5$\\
       calfs & $w_q = 0.3$\\
       feet & $w_q = 0.2, w_p = 0.1$\\
       target object & $w_q = 1, w_p = 1$ \\
        \bottomrule
    \end{tabular}
    \caption{Unnormalized weights of each body link of the character and the target object when computing the reward regarding body pose errors (cf. $r_\text{body}$ in Equation~\ref{eq:r_body}). All weights unlisted are zero.
    The weights are normalized such that $\sum_b w_{p,b}=1$ and $\sum_b w_{q,b}=1$ before computing $r_\text{body}$.}
    \label{tab:rl_weights}
\end{table}

\begin{table}[t]
    \centering
\begin{tabular}{lc}
    \toprule
    \textbf{Parameter} & \textbf{Value}\\
    \midrule
    policy network learning rate & $5 \times 10^{-6}$\\
    critic network learning rate & $1 \times 10^{-4}$\\
    reward discount factor ($\gamma$) & $0.95$ \\
    GAE discount factor ($\lambda$) & $0.95$ \\
    surrogate clip range ($\epsilon$) & $0.2$ \\
    number of PPO workers & \\
    \quad (simulation environment instances) & $2048$ \\
    PPO replay buffer size & $2048\times8$ \\
    PPO batch size & $256$ \\
    PPO optimization epochs & $5$ \\
  \bottomrule
  \end{tabular}
    \caption{Hyperparameters of policy training for physics tracker.}
    \label{tab:rl_hyper}
    
\end{table}


For full-body pose evaluation,
we consider the orientation and position accuracy of the character's body links and the target object:
\begin{equation}\label{eq:r_body}\begin{split}
    r_\text{body} = & 0.5\exp\left(-15 \sum_{b\in\mathcal{B}} w_{q,b} ||\mathbf{e}_{q,b}||^2\right) \\
    & \quad + 0.5\exp\left(-15 \sum_{b\in\mathcal{B}} w_{p,b} ||\mathbf{e}_{p,b}||^2\right),
\end{split}
\end{equation}
where $\mathcal{B}$ is the set of all body links of the character and the activated target object at the current time step, $\mathbf{e}_{q,b}$ is the orientation error between each body link $b$ of the simulated character and that in the tracking motion, $\mathbf{e}_{p,b}$ is the position error, and $w_{q,b}$ and $w_{p,b}$ are the associated weights.
We measure $\mathbf{e}_{q,b}$ using the radian of the angle needed to rotate the link $b$ to the target orientation,
while $\mathbf{e}_{p,b}$ is measured by the Euclidean distance from the current position of link $b$ to the target position.
All weight values are listed in Tab.~\ref{tab:rl_weights}. For position tracking, we focus only on the position of the root link and end effectors to ensure that the trajectory of end effectors (foot and hand excluding fingers) matches the tracking target,
while relying on orientation tracking to ensure the naturalness of full-body poses.

For hand pose evaluation,
to ensure that the object would be manipulated by hands correctly,
we consider the fingers' relative position to the target manipulating object when the hand is close to the object, or their relative position to the wrist when the hand is far away from the target object:
\begin{equation}
    r_\text{hand} = \exp\left(- \frac{5}{|\mathcal{F}|}\sum_{f\in\mathcal{F}} \alpha ||\mathbf{e}_{f,o}|| + (1-\alpha)  ||\mathbf{e}_{f,w}||\right),
\end{equation}
where $\mathcal{F}$ is the set of fingers, $\mathbf{e}_{f,o}$ is the position error between each finger $f$ of the simulated character and the tracking target position in the object's local system, $\mathbf{e}_{f,w}$ is the position error in the wrist's local system, and $\alpha$ is an interpolation coefficient depending on the distance between the hand and object in the generated tracking motions. We define $\alpha=1$ when the hand is equal to or less than 0.25m from the target object and $\alpha=0$ if the hand is at least 1m far away from the object in the tracking motion.


$r_\text{energy}$ takes into account the character's end effectors to prevent jittering or abrupt movement with larger acceleration:
\begin{equation}
    r_\text{energy} = \exp\left(-\frac{1}{900}\sum_{e\in\mathcal{E}} ||\mathbf{a}_e||^2\right),
\end{equation}
where $\mathcal{E}$ is the set of the character's key end effectors (feet and hands excluding fingers), and $\mathbf{a}_e$ is the linear acceleration of the end effector $e$.

In contrast to recent works~\cite{brahmbhatt2019contactgrasp,wang2023physhoi,luo2024grasping} on object manipulation, our reward function does not rely on the contact states between the hands (or fingers) and objects. This is because the contact data derived from the kinematic motions generated by diffusion models can be unreliable due to artifacts (cf.~\cref{fig:physics_qualitative_results}). Moreover, this approach aligns with the common scenario in which detailed contact information is typically absent in motions obtained from motion capture and video demonstrations~\cite{pollard2005physically}.


\subsection{Importance Sampling Strategy}

Given that the generated kinematic motions can span several minutes and involve interactions with multiple objects, it would be highly inefficient to perform motion tracking from the very beginning of the sequence. Since locomotion without object manipulation and the phase following object grasping are relatively easier to learn, we adopt an importance sampling strategy that focuses on pre-grasp states to facilitate policy training.
Specifically, we partition the parallel simulated physics environments into $|\mathcal{O}|$ batches, with each batch assigned to a specific manipulating object. In each batch, simulated characters are initialized in a random pre-grasp pose while interacting with the assigned object. This approach increases the likelihood that the policy learns the transition from the pre-grasp to the grasping phase, and subsequently, the locomotion that follows a successful grasp.
Moreover, this strategy ensures that the policy consistently learns to interact with multiple objects simultaneously, thereby avoiding potential local minima during policy exploration.




\section{Inference Speed}
Achieving real-time performance remains an open challenge for both diffusion-based and optimization-based methods. In our low-level motion generator, producing a 4-second motion sequence takes about four seconds for each diffusion model in stages 1, 3, and 4, while the optimization process in stage 2 requires roughly three minutes. Additionally, training the physics tracker on a 4-second interaction segment takes approximately 2 hours.
All experiments were run on a single NVIDIA RTX 4000 GPU.

\begin{figure}[t]
    \centering
     % \captionsetup{aboveskip=0pt, belowskip=0pt}
\includegraphics[width=0.48\textwidth]{figures/8_physhoi.pdf}
% \vspace{-7mm}
    \caption{PhysHOI~\cite{wang2023physhoi} fails to track kinematic motion, while our policy successfully drives the character to interact with diverse objects.}
    \label{fig:physhoi}
    % \vspace{-5mm}
\end{figure}

\section{Additional Experiemental Analysis}


\subsection{Implementation Details of Baseline}
Since no prior work addresses the task of generating full-body motion, finger motion, and object motion from language instructions, we adapt prior works CHOIS~\cite{li2023controllable} and GRIP~\cite{taheri2023grip} as fair baselines to compare against our system. CHOIS generates full-body motion and object motion from text and initial states, which also serves as our Stage 1. The output of CHOIS will serve as input to GRIP. GRIP consists of an arm denoising model and a two-stage finger motion generation model. It takes full-body motion (without finger motion) and object motion as input to denoise arm motion. It then proposes a hand sensor to extract hand-object spatial features and predict initial finger poses. In the second stage of finger motion generation, GRIP leverages these predicted finger poses to compute proximity features, which serve as input to refine the finger motion.

\subsection{Results of Navigation Module}
We report the evaluation scores for the navigation module in \cref{tab:navigation}, using the metrics introduced in~\cref{subsec:low_level_eval}. The module is evaluated on the test split of the HumanML3D dataset \cite{guo2022generating}. It demonstrates that the module closely adheres to the input waypoints, with a relatively small joint position error when compared to the ground truth data, while maintaining reasonable foot sliding and foot height.

\begin{table}[t]
\centering
% \small
\begin{adjustbox}{max width=0.5\textwidth}
\begin{tabular}{lcccccccccccccccccc}
\toprule
& $T_{xy}$ & $H_{\text{feet}}$ & $FS$ & $\text{MPJPE}$  \\
\midrule

Ours
& 4.00 & 1.48 & 0.58 & 10.87 \\
\bottomrule
\end{tabular}
\end{adjustbox}

\caption{Navigation synthesis on the HumanML3D dataset \cite{guo2022generating}. The joint error is relatively small, demonstrating that the model aligns well with the input waypoints.}
\label{tab:navigation}
\end{table}



\subsection{Results of Physics Tracker}
We compare our physics tracker against PhysHOI~\cite{wang2023physhoi}, which trains reinforcement learning policies to control a humanoid character playing basketball. 
For our task, we have PhysHOI track the kinematic motion generated by our low-level motion generator. As shown in \cref{fig:physhoi}, PhysHOI fails to grasp objects, whereas our method successfully completes the task.

\subsection{Human Perception Study Details}
We show the interface of the human perceptual study for the high-level planner in Figure~\ref{fig:llm_human_study} and the low-level motion generator in Figure~\ref{fig:motion_human_study}. 



% \section{Evaluation of High-Level Planner}

 
% % We show an example for processing a task using our LLM planner in Figure~\ref{fig:llm_our_example}.

% \begin{figure}[!ht]
% \begin{center}
%     \includegraphics[width=0.9\columnwidth]{figures/tasks15_llm_eval.pdf}
%     % \vspace{-3mm}
%     \caption{Tasks for evaluating target layout generation.}
%     \label{fig:llm_all_tasks}
%     \vspace{-5mm}
% \end{center}
% \end{figure}

% \begin{figure*}[!ht]
%     \centering
%      % \captionsetup{aboveskip=0pt, belowskip=0pt}
%      \vspace{-4mm}
% \includegraphics[width=0.9\textwidth]{figures/baseline_prompt_new.pdf}
% \vspace{-4mm}
%     \caption{Prompts for the baseline of target layout generation evaluation.} 
%     \label{fig:llm_baseline}
% \end{figure*}




% We show the prompt for the baseline in Figure~\ref{fig:llm_baseline}.
% We provide descriptions of 15 tasks used for evaluating target layout generation in Figure~\ref{fig:llm_all_tasks}. 


\begin{figure*}[h]
    \centering
     % \captionsetup{aboveskip=0pt, belowskip=0pt}
\includegraphics[width=\textwidth]{figures/llm_human_study.pdf}
    \caption{Interface of human perceptual study for the high-level planner.} 
    \label{fig:llm_human_study}
    \vspace{-2mm}
\end{figure*}

\begin{figure*}[h]
    \centering
     % \captionsetup{aboveskip=0pt, belowskip=0pt}
\includegraphics[width=\textwidth]{figures/motion_human_study.pdf}
    \caption{Interface of human perceptual study for the low-level motion generator.} 
    \label{fig:motion_human_study}
    \vspace{-2mm}
\end{figure*}

% Bibliography
% \clearpage
% {
%     \small
%     \bibliographystyle{ieeenat_fullname}
%     \bibliography{sample-bibliography}
% }


% \end{document}


\end{document}
